\documentclass{situationist-preprint}

\title{Classical Identity and the Counting Problem}
\msid{REL-13-b}
\msfamily{REL-13}
\msdiscipline{Philosophy of Religion}
\mstype{Companion manuscript}
\mskeywords{Trinity; counting; absolute identity; relative identity; tritheism; numerical sameness}

\begin{document}
\maketitle

\begin{abstract}
REL-13-a proves that the doctrine of the Trinity is consistent in any logic of relative identity extending Geach's system. A defender of classical identity may reply that the relative-identity reading changes the subject, and that the logical problem should be solved with identity as it is. We show that the counting problem leaves no such option. In any language that contains classical identity, if the relation by which the persons are said to be one God is identity, the doctrine is inconsistent. Every consistent reading must count Gods by an equivalence relation other than identity, which is not substitutive. This is the relative-identity reading, and it is available to one who holds that identity is absolute. The social, Latin and constitution models of the Trinity each confirm the dilemma by taking one of its options.
\end{abstract}

\section{The counting problem}

The Athanasian Creed says that the Father is God, the Son is God and the Spirit is God, and that there are not three Gods but one God. The Father is not the Son, the Son is not the Spirit, and the Father is not the Spirit. Read with classical identity and standard counting, the three are distinct objects each of which is a God, so there are at least three Gods, and there is exactly one \citep{Cartwright1987}. This is the counting problem.

We use the claims (T1) to (T6) of REL-13-a's Definition~2.6, with $D$ for \enquote{God}, $P$ for \enquote{person}, $f$, $s$, $h$ for the three persons, and $=_D$, $=_P$ for the relations of being the same God and the same person. We now add classical identity, $=$, governed by reflexivity and the substitutivity of identicals for every formula.

\begin{proposition}\label{p:identity}
If $=_D$ is identity, then (T2), (T3) and (T4) are inconsistent, given that $=_P$ is reflexive on persons.
\end{proposition}

\begin{proof}
By (T3), $f =_D s$, so $f = s$. By (T2) and the reflexivity of $=_P$ on persons, $f =_P f$. Substituting $s$ for the second occurrence of $f$, by substitutivity of identity, gives $f =_P s$, contrary to (T4).
\end{proof}

\begin{proposition}\label{p:count}
If Gods are counted by identity, that is, if (T5) is read as \enquote{there is an $x$ such that $Dx$ and every $y$ such that $Dy$ is identical to $x$}, then (T1) and (T5) together entail $f = s = h$, and with (T4) and the reflexivity of $=_P$ on persons they are inconsistent.
\end{proposition}

\begin{proof}
By (T1), $f$, $s$ and $h$ are Gods; by (T5) so read, each is identical to the one God $x$, so $f = s = h$. The argument of Proposition~\ref{p:identity} then derives $f =_P s$.
\end{proof}

Propositions~\ref{p:identity} and~\ref{p:count} do not depend on the persons' being counted by identity; they need only that \enquote{same person} is reflexive on persons, which is the first principle of Geach's system (REL-13-a, Definition~2.4). They locate the problem precisely: what cannot be identity, on pain of inconsistency, is the relation by which the persons are one God and by which Gods are counted.

\begin{corollary}\label{c:not}
In every consistent reading of (T1) to (T6) in a language with classical identity, \enquote{same God} is an equivalence relation on Gods that is not identity and is not substitutive for the predicate \enquote{is the same person as the Father}.
\end{corollary}

\begin{proof}
By Proposition~\ref{p:identity}, $=_D$ is not identity. By (T2), (T3) and (T4), $f =_D s$, $f$ satisfies \enquote{$\cdot =_P f$} and $s$ does not, which is the failure of substitutivity of REL-13-a's Proposition~4.5. That $=_D$ is an equivalence relation on Gods is the content of Geach's system for the sortal $D$.
\end{proof}

Corollary~\ref{c:not} is what the classical theorist must accept. It does not require denying that identity is absolute, as LOG-12 concludes it is. It requires that the creedal \enquote{one God} be read as a count by \enquote{same God}, and that \enquote{same God} not be identity. REL-13-a's Remark~4.6 records that its consistency theorem survives the addition of absolute identity: in the extended model, the persons are three objects that are one God. Classical identity and the doctrine are compatible exactly when Gods are counted by \enquote{same God}.

\section{The classical models}

The main models of the Trinity that keep classical identity confirm the dilemma; each takes one of its options.

\emph{Social models} \citep{Swinburne1994,Hasker2013} hold that the persons are three distinct divine individuals, and read \enquote{one God} as a claim about their unity of will, nature or community. If \enquote{God} in (T1) is predicated of the persons, Gods are counted by a relation other than identity, the unity of the community, and the model is the relative-identity reading with a particular interpretation of $=_D$. If instead \enquote{one God} names the community, (T1) is read as \enquote{is a member of God}, and the doctrine's claim that each person is God is given up in its literal form.

\emph{Latin models} \citep{Leftow2004} hold that there is one God, who lives three lives, and that the persons are God as living each of them. The persons are then not three objects distinct by identity; their distinctness in (T4) is a distinction between lives, or between God under three descriptions. The model keeps (T5) read by identity and gives up (T4) read by identity, as Proposition~\ref{p:identity} requires.

\emph{Constitution models} \citep{BrowerRea2005} hold that the persons are numerically the same God without being identical, as a statue and its bronze are numerically the same material object without being identical. Numerical sameness without identity is an equivalence relation, weaker than identity, by which the persons are counted as one God. The model is the relative-identity reading, supplied with a metaphysical interpretation of $=_D$ drawn from the theory of material constitution.

\begin{objection}
Counting by a relation other than identity is not counting. If there are three objects each of which is a God, there are three Gods, whatever relation they bear to one another; the relative-identity reading is tritheism stated in other words.
\end{objection}

\begin{reply}
Counting is always counting of things under a sortal, and how many there are depends on the criterion of sameness the sortal supplies \citep{Geach1962}. Asked how many words are on a page, one may count tokens or types and get different numbers, and neither answer is tritheism about words. The claim that there are three Gods is false on the reading of REL-13-a, since every God is the same God as the Father. What the reading accepts, when absolute identity is added, is that there are three objects each of which is a God; whether this is objectionable is the question whether Gods should be counted as objects are, which is not a question of logic. Proposition~\ref{p:count} shows that, if they are, no reading of the doctrine is consistent.
\end{reply}

\section{Conclusion}

The counting problem is not solved by keeping classical identity and adjusting the theology. In any consistent reading of the doctrine, the persons are one God by a relation that is not identity, and Gods are counted by that relation. This is the relative-identity reading whose consistency REL-13-a proves, and it is compatible with the absoluteness of identity.

\begin{thebibliography}{99}\small
\bibitem[Brower and Rea(2005)]{BrowerRea2005} Brower, J.~E. and Rea, M.~C. (2005). Material constitution and the Trinity. \emph{Faith and Philosophy}, 22(1), 57--76.
\bibitem[Cartwright(1987)]{Cartwright1987} Cartwright, R. (1987). On the logical problem of the Trinity. In \emph{Philosophical Essays}, MIT Press.
\bibitem[Geach(1962)]{Geach1962} Geach, P.~T. (1962). \emph{Reference and Generality}. Cornell University Press.
\bibitem[Hasker(2013)]{Hasker2013} Hasker, W. (2013). \emph{Metaphysics and the Tri-Personal God}. Oxford University Press.
\bibitem[Leftow(2004)]{Leftow2004} Leftow, B. (2004). A Latin Trinity. \emph{Faith and Philosophy}, 21(3), 304--333.
\bibitem[Swinburne(1994)]{Swinburne1994} Swinburne, R. (1994). \emph{The Christian God}. Oxford University Press.
\end{thebibliography}

\end{document}
